\documentclass[10pt,a4paper]{amsart}
\usepackage{amsmath,amssymb,amsthm,graphicx}
% For commutative diagrams we use Xy-pic
\usepackage[all]{xy}

% We use 2cell for 2-commutative diagrams.
\xyoption{2cell}
% Theorem environments.
%
\theoremstyle{plain}
\newtheorem{theorem}[subsection]{Theorem}
\newtheorem{proposition}[subsection]{Proposition}
\newtheorem{lemma}[subsection]{Lemma}

\theoremstyle{definition}
\newtheorem{definition}[subsection]{Definition}
\newtheorem{example}[subsection]{Example}
\newtheorem{exercise}[subsection]{Exercise}
\newtheorem{situation}[subsection]{Situation}

\theoremstyle{remark}
\newtheorem{remark}[subsection]{Remark}
\newtheorem{remarks}[subsection]{Remarks}

\numberwithin{equation}{subsection}

% Macros
%
\def\lim{\mathop{\mathrm{lim}}\nolimits}
\def\colim{\mathop{\mathrm{colim}}\nolimits}
\def\Spec{\mathop{\mathrm{Spec}}}
\def\Hom{\mathop{\mathrm{Hom}}\nolimits}
\def\Ext{\mathop{\mathrm{Ext}}\nolimits}
\def\SheafHom{\mathop{\mathcal{H}\!\mathit{om}}\nolimits}
\def\SheafExt{\mathop{\mathcal{E}\!\mathit{xt}}\nolimits}
\def\Sch{\mathit{Sch}}
\def\Mor{\mathop{\mathrm{Mor}}\nolimits}
\def\Ob{\mathop{\mathrm{Ob}}\nolimits}
\def\Sh{\mathop{\mathit{Sh}}\nolimits}
\def\NL{\mathop{N\!L}\nolimits}
\def\CH{\mathop{\mathrm{CH}}\nolimits}
\def\proetale{{pro\text{-}\acute{e}tale}}
\def\etale{{\acute{e}tale}}
\def\QCoh{\mathit{QCoh}}
\def\Ker{\mathop{\mathrm{Ker}}}
\def\Im{\mathop{\mathrm{Im}}}
\def\Coker{\mathop{\mathrm{Coker}}}
\def\Coim{\mathop{\mathrm{Coim}}}

% Boxtimes
%
\DeclareMathSymbol{\boxtimes}{\mathbin}{AMSa}{"02}

%
% Macros for moduli stacks/spaces
%
\def\QCohstack{\mathcal{QC}\!\mathit{oh}}
\def\Cohstack{\mathcal{C}\!\mathit{oh}}
\def\Spacesstack{\mathcal{S}\!\mathit{paces}}
\def\Quotfunctor{\mathrm{Quot}}
\def\Hilbfunctor{\mathrm{Hilb}}
\def\Curvesstack{\mathcal{C}\!\mathit{urves}}
\def\Polarizedstack{\mathcal{P}\!\mathit{olarized}}
\def\Complexesstack{\mathcal{C}\!\mathit{omplexes}}
% \Pic is the operator that assigns to X its picard group, usage \Pic(X)
% \Picardstack_{X/B} denotes the Picard stack of X over B
% \Picardfunctor_{X/B} denotes the Picard functor of X over B
\def\Pic{\mathop{\mathrm{Pic}}\nolimits}
\def\Picardstack{\mathcal{P}\!\mathit{ic}}
\def\Picardfunctor{\mathrm{Pic}}
\def\Deformationcategory{\mathcal{D}\!\mathit{ef}}
\usepackage{cleveref}
\crefname{originaldefinition}{original definition}{original definitions}
\Crefname{originaldefinition}{Original definition}{Original definitions}
\makeatletter
\newlabel{algebra-definition-standard-smooth}{{00T6}{source}{}{sourcepointer-1}{}}
\newlabel{algebra-definition-standard-smooth@cref}{{[originaldefinition][0][]00T6}{[1][0][]source}{}{}{}}
\makeatother
\title{A frobenius linear endomorphism of a finite dimensional vector space over an algebraically closed field has surjective F minus one}
\author{The Stacks Project Authors}
\date{stacks-project:etale-cohomology.tex:local-snapshot-20261011; GNU FDL 1.2 or later}
\begin{document}
\maketitle
\noindent\textit{Editorial note.} All delivered mathematical text is the author's own, copied from the pinned Stacks Project chapter file etale-cohomology.tex; no statement or proof step was rewritten, reordered or removed.\par\medskip
\noindent\textit{Editorial note.} The proof is the author's complete argument: the reduction, the matrix computation, the construction of A and the point count are all inside the counted span.\par\medskip
\noindent\textit{Editorial note.} The label lemma-F-1 occurs in no other staged Stacks chapter, so no cross-chapter restatement had to be weighed for this card.\par\medskip
\noindent\textit{Editorial note.} The single labelled reference occurring inside the counted spans points to a chapter of the Project that is not part of this card's source tree; it is carried in assembly.external\_references as a non-load-bearing navigation pointer recorded at the invoking line.\par\medskip
\noindent\textit{Editorial note.} The preamble is the chapter's own native preamble excerpt already used for this source: the Xy-pic setup (native preamble lines 14-18) and the native theorem declarations with the mathematical macros (native preamble lines 154-217).\par\medskip
\noindent\textit{Allowed standard prerequisites (editorial).}\begin{itemize}
\item Frobenius linear map on a vector space over a field of characteristic p. Same chapter, the notion of an additive map F with F(lambda v) = lambda\textasciicircum{}p F(v) for lambda in k. Supplies the structure on V that makes the equation F - 1 meaningful and the algebra A below finite over k.
\item kernels and graded pieces of a filtration by F-stable subspaces. Same chapter, the standard reduction argument allowing a statement to be checked on graded pieces of a filtration by F-stable subspaces. Supplies the reduction to the case where the kernel of F is zero.
\item Standard smooth algebras and counting the points of their spectra. Algebra, the definition of a standard smooth algebra together with the fact that a standard smooth algebra over an algebraically closed field has as many points as its k-dimension. Supplies the conclusion that Spec(A) has p\textasciicircum{}n points, hence surjectivity of F - 1.
\end{itemize}\medskip
\ifdefined\hypertarget\hypertarget{sourcepointer-1}{}\fi
\noindent\textit{Source reference.} Original definition 00T6: sources/stacks-etale-cohomology/pinned-original-source.tex, source lines 9682--9682. Cross-chapter navigation pointer occurring in the proof at source line 9682, where the author records that the displayed algebra is standard smooth and names the definition in the chapter Algebra. The target chapter is not part of this card's source tree, so the pointer records the invoking line; it only resolves the reference for the reader, and no bibliography citation is involved.\par\medskip
\begin{lemma}
\label{lemma-F-1}
Let $k$ be an algebraically closed field of characteristic $p > 0$.
Let $V$ be a finite dimensional $k$-vector space. Let $F : V \to V$
be a frobenius linear map, i.e., an additive map such that
$F(\lambda v) = \lambda^p F(v)$ for all $\lambda \in k$ and $v \in V$.
Then $F - 1 : V \to V$ is surjective with kernel a finite dimensional
$\mathbf{F}_p$-vector space of dimension $\leq \dim_k(V)$.
\end{lemma}

\begin{proof}
If $F = 0$, then the statement holds. If we have a filtration of $V$ by
$F$-stable subvector spaces such that the statement holds for each
graded piece, then it holds for $(V, F)$. Combining these two remarks
we may assume the kernel of $F$ is zero.

\medskip\noindent
Choose a basis $v_1, \ldots, v_n$ of $V$ and write
$F(v_i) = \sum a_{ij} v_j$. Observe that $v = \sum \lambda_i v_i$
is in the kernel if and only if $\sum \lambda_i^p a_{ij} v_j = 0$.
Since $k$ is algebraically closed this implies the matrix $(a_{ij})$
is invertible. Let $(b_{ij})$ be its inverse. Then to see that $F - 1$
is surjective we pick $w = \sum \mu_i v_i \in V$ and we try to solve
$$
(F - 1)(\sum \lambda_iv_i) =
\sum \lambda_i^p a_{ij} v_j - \sum \lambda_j v_j = \sum \mu_j v_j
$$
This is equivalent to
$$
\sum \lambda_j^p v_j - \sum b_{ij} \lambda_i v_j = \sum b_{ij} \mu_i v_j
$$
in other words
$$
\lambda_j^p - \sum b_{ij} \lambda_i = \sum b_{ij} \mu_i,
\quad j = 1, \ldots, \dim(V).
$$
The algebra
$$
A = k[x_1, \ldots, x_n]/
(x_j^p - \sum b_{ij} x_i - \sum b_{ij} \mu_i)
$$
is standard smooth over $k$
(Algebra, Definition \ref{algebra-definition-standard-smooth})
because the matrix $(b_{ij})$ is invertible and the partial derivatives
of $x_j^p$ are zero. A basis of $A$ over $k$ is the set of monomials
$x_1^{e_1} \ldots x_n^{e_n}$ with $e_i < p$, hence $\dim_k(A) = p^n$.
Since $k$ is algebraically closed we see that $\Spec(A)$ has exactly
$p^n$ points. It follows that $F - 1$ is surjective and every fibre
has $p^n$ points, i.e., the kernel of $F - 1$ is a group with $p^n$ elements.
\end{proof}

\end{document}
